\section{Parity-separated roots: thirteen integers, degree 156}
\label{sec:parity-three-root13}
This endpoint uses thirteen unrestricted integers, has exact degree $156$, and
has a division-free evaluator of $127$ binary operations: $75$ multiplications,
$27$ additions and $25$ subtractions. It eliminates both the strong Pell root and
the triangular positivity coordinate from the fixed fifteen-variable free-scale
source. Its zero existence is equivalent to that source by maps on every zero,
not just a canonical subclass. The same ordinary source34 compiler/calibration
therefore supplies the ZFC-provable equivalence between zero nonexistence and
syntactic $\mathrm{Con}(\mathrm{ZFC})$; this is not a proof of consistency or a new
end-to-end formal-kernel certificate.

\subsection{Complete polynomial}
The unknowns are
\[
 j,C,z,h,r,v,w,T,c,n,y,p_1,p_2\in\mathbb Z.
\]
The six constants $\beta,K_0,M_C,M_F,A_c,B_c$ are the unchanged complete source
payloads, with $M_F$ already shifted. Define only polynomial abbreviations:
\begin{gather*}
 J=6j+1,\quad q_m=\beta J,\quad q=q_m+1,\quad W=12w+4,\\
 X=qW,\quad Y=v,\quad F=(K_0+W)C-zq_m,\\
 R=(q(q-F)-C)(q^2-1)+(M_C+qM_F)J,\\
 E=XY,\quad a=E+Y,\quad H=4a+3,\quad\Delta=a^2+H,\\
 k=R+1+hE,\quad \eta=c-Yk,\quad L=E(Yk),\\
 I=(2r+1)c^2,\quad D_s=\Delta I^2,\quad F_s=1+D_s,\\
 Q=\Delta D_s,\quad V=F_s(cT-R)-c,\\
 M=2q^3,\quad U=MH,\quad K_{\rm lo}=X(2X-3c)-2,\\
 Z=Un+MK_{\rm lo}+H(Y-q^3),\\
 D_{\rm err}=(A_cF-q)^2+(B_cC-q)^2,\\
 B_{\rm ud}=(Yk)^2(q^2-D_{\rm err})
       -q^2\bigl(4(Y\eta)^2+(kq^3)^2\bigr),\\
 G=4(B_{\rm ud}-p_1^2-p_2^2)-3,\qquad D_t=1+L(L+k),\\
 C_1=(Z+Uk)^2+4U(LZ-U),\qquad
 C_2=V^2-(Q-1)y(2V+y)-1,\\
 g_a=U^2G,\quad d_b=4U^2D_t,\quad f_c=16U^2F_s,\quad g_c=g_a+f_c,\\
 J_N=(C_1-g_c)^2-4(g_af_c+d_bg_c),\\
 N=J_N^2-64g_ad_bf_c(C_1+d_b).
\end{gather*}
The single equation is
\[
 \boxed{P=(N+1)(C_2+1)-1=0.}
\]
Sharing $D_s$ avoids recomputing $F_s-1$ and saves one subtraction. All powers
are fixed natural powers; the literal DAG has no divisions or radicals.

\subsection{Norm and exact conjunction}
For proof only, put $b=2L+k$, $R_N=Z+Ub$ and $\rho=R_N/U$.
The first-carrier identity gives $R_N^2=C_1+d_b$. The symmetric norm
\[
 F_8(S;g,d,f)=\bigl[(S-g-d-f)^2-4(gd+gf+df)\bigr]^2-64gdfS
\]
satisfies the polynomial identity
\[
 F_8(x^2;g,d,f)=\prod_{\epsilon\in\{\pm1\}^3}
   \bigl(x-\epsilon_1\sqrt g-\epsilon_2\sqrt d-\epsilon_3\sqrt f\bigr).
\]
Substitution $S=C_1+d_b$, $(g,d,f)=(g_a,d_b,f_c)$ gives exactly $N$.
Equivalently,
\[
 \frac{N}{U^8}=\prod_{\epsilon\in\{\pm1\}^3}
 \bigl(\rho-\epsilon_1\sqrt G-2\epsilon_2\sqrt{D_t}
                         -4\epsilon_3\sqrt{F_s}\bigr). \tag{T1}
\]
Since $N$ is a square modulo $4$, $N\ne-2$. The integer product
$(N+1)(C_2+1)=1$ therefore forces both factors to be $1$; it is equivalent to
$N=C_2=0$ with no sign premise on either residual. This is a tailored use of the
standard unit-product device, not a claimed new general conjunction theorem.

\subsection{Unconditional signs and square-class separation}
The authentic odd $\beta>1$ makes $q$ even and nonzero. Also $H$ is odd and
nonzero, so $U\ne0$. Since $|W|\ge4$, $X$ is even with $|X|\ge8$.
Thus $a=(X+1)Y$ is either $0$ or has $|a|\ge7$, giving
\[
 \Delta=(a+2)^2-1\ge3,\qquad F_s=1+\Delta I^2\ge1.
\]
Writing $D=XY^2$, we have even $D$ and $D(D+1)\ge0$, whence
$D_t=1+D(D+1)k^2$ is a positive odd integer. Also $G\equiv1\pmod4$.
If $G<0$, every factor in (T1) has a nonzero imaginary part from its sole
imaginary radical, while $\rho,\sqrt{D_t},\sqrt{F_s}$ are real. Hence at a zero
$G>0$; $G=0$ is excluded by its residue.

Write the positive radicands uniquely as
\[
 G=u_1^2d_1,\qquad D_t=u_2^2d_2,\qquad F_s=u_3^2d_3,
\]
where the $d_i$ are positive squarefree and the $u_i$ are positive integers.
The numbers $u_1,u_2,d_1,d_2$ are odd. A vanishing norm factor gives
\[
 \rho=\epsilon_1u_1\sqrt{d_1}
       +2\epsilon_2u_2\sqrt{d_2}+4\epsilon_3u_3\sqrt{d_3}\in\mathbb Q.
\]
Distinct squarefree radicals, including $\sqrt1$, are linearly independent over
$\mathbb Q$: sign-change automorphisms in the finite multiquadratic extension,
averaged against a chosen character, isolate each coefficient.
Group terms with equal squarefree part. A group containing the $G$ term has an
odd coefficient sum. A group containing $D_t$ but not $G$ has coefficient sum
$2\pmod4$. A remaining group contains only the nonzero $F_s$ term. Thus no
irrational square-class group can cancel. An even squarefree class, in
particular, could contain only $F_s$ and is also impossible.
Consequently all three radicals are integers. With the signs of the vanishing
factor, write
\[
 g=\epsilon_1\sqrt G,\qquad \tau=\epsilon_2\sqrt{D_t},\qquad
 f=\epsilon_3\sqrt{F_s};\qquad \rho=g+2\tau+4f.
\]
Here $g$ is odd; no oddness of $F_s$ was assumed. This is an exact valuation
argument within common square classes, not an inference from pairwise distinct
radicands or finite tests.

\subsection{Complete maps and canonical completeness}
From any new zero restore the old coordinates by
\[
 f_{\rm old}=f,\qquad p_{3,\rm old}=(g-1)/2,\qquad
 n_{\rm old}=n-g-4f,
\]
retaining every other coordinate. All three restored coordinates are integral.
Affineness of $Z$ gives $Z_{\rm old}/U=2\tau-b$, making its first carrier zero.
The identity $f^2=F_s=1+\Delta I^2$ makes its strong row zero; its auxiliary row
is unchanged because the new $F_s$ equals the old $f^2$. Finally,
\[
 p_{3,\rm old}(p_{3,\rm old}+1)=(G-1)/4,
\]
so its guard is precisely
$B_{\rm ud}=1+p_1^2+p_2^2+p_{3,\rm old}(p_{3,\rm old}+1)$.
This is a literal source zero, after which its ordinary soundness theorem
applies. No preliminary CRT, positive orientation or canonical restriction was
used.

Conversely, from any old zero delete $f,p_3$ and put
\[
 n_{\rm new}=n_{\rm old}+(2p_{3,\rm old}+1)+4f_{\rm old}.
\]
The old rows imply $F_s=f_{\rm old}^2$, $G=(2p_{3,\rm old}+1)^2$, and
$Z_{\rm old}/U+b=2\tau$ with integral $\tau^2=D_t$. The latter integrality is
explicitly proved by the pinned source's monic carrier argument. One factor in
(T1) is therefore zero and the auxiliary row is unchanged. This map applies to
every old zero, in particular to the exact canonical source34 family supplied
from a finite refutation. No extra canonical tuning or simultaneous
representation theorem is required.

\subsection{Exact coefficient and verification}
The degrees of $C_1,F_s,D_t,G,Q,V,C_2,N$ are respectively
$22,12,18,18,18,17,36,120$. Put
\begin{gather*}
 q_0=6\beta j,\quad X_0=12q_0w,\quad Y_0=v,\quad
 R_0=q_0^4z-12q_0^3wC,\\
 a_0=X_0Y_0,\quad U_0=8q_0^3X_0Y_0.
\end{gather*}
The leading forms are
\begin{gather*}
 N_{120}=2^8U_0^8R_0^8q_0^{16}(q_0^8+4X_0^2Y_0^4)^2,\\
 (C_2)_{36}=32a_0^6r^4c^8R_0y.
\end{gather*}
Their product is the unique degree-$156$ part, and
\[
 [j^{106}w^{14}v^{14}z^9r^4c^8y]P
       =2^{171}3^{120}\beta^{106}>0.
\]
The exponent sum is $156$, so this is actual degree attainment on the fixed
coefficient slice. All thirteen variables and six coefficient ports are live;
the unchanged numerals occupy $615883$ uncompressed bytes.

The independent normal and optimized audits authenticate the source and numeral
bytes, check every named register, multiply all eight signed-root factors,
verify the centered norm schedule and lift, prove the coefficient and degree,
check literal cost and liveness, and reject all live-gate mutations. Finite
checks are not substituted for the square-class proof or the unchanged
unbounded source34 compiler/refutation theorem. This endpoint is a new
variable/degree/cost tradeoff, not a domination claim over the fourteen-variable
degree-$88$ or fifteen-variable degree-$44$ constructions.
